\section{Evaluation}
\label{sec:evaluation}

The evaluation combines an address-level production evidence audit, controlled
implementation checks, generalization and boundary sweeps of those checks, an
independent from-scratch reimplementation, and a public-testnet enforcement
execution. The first pins the active Gnosis committee and computes both its
supported floor and the evidence needed to change it. Controlled profiles
exercise the certificate branches and check the reference implementation; a
further sweep repeats them at other committee shapes and parameter extremes,
and a separate reimplementation, built only from this paper's own formulas,
cross-checks the numbers below. The Chiado pilot tests a machine-readable
evidence-to-penalty transition.

We ask what the production record certifies; how activation order and
acquisition mechanism change fixed-resistance costs; whether local defensive
measures miss coordinated covers; whether the four-of-seven findings survive
other committee shapes and parameter extremes; and whether a public testnet
penalty can execute without being transferred to production.

The controlled checks use deterministic Python instances; random lattice
checks use seed \(20260714\). All ten recorded timing runs are retained. The
scaling machine ran Python~3.11.1 on an Intel Core i7-13650HX with 16.85~GB
physical memory. The enforcement implementation uses Rolling Shutter~v1.4.4,
Solidity~0.8.28, Hardhat, Node.js~22, and Gnosis Chiado. Seven validator and
seven Keyper processes run on one host. The controlled environments check
implementation paths and complexity behavior; the production audit supplies
real committee identities but does not infer unobserved member resistance.

\subsection{Controlled Checks}
\label{subsec:evaluation-setup}
\label{subsec:certificate-semantics-evaluation}
\label{subsec:activation-ladder-evaluation}
\label{subsec:defensive-allocation-evaluation}
\label{subsec:parameter-sensitivity-evaluation}
\label{subsec:scalability-evaluation}
\label{subsec:evaluation-baselines}

\begin{table}[!b]
\centering
\captionsetup{font=scriptsize,skip=2pt}
\fontsize{7}{7.5}\selectfont
\setlength{\tabcolsep}{2pt}
\renewcommand{\arraystretch}{0.92}
\caption{Controlled checks; positive inputs are non-production fixtures.}
\label{tab:controlled-results}
\label{tab:evaluation-layers}
\label{tab:certificate-branch-results}
\label{tab:activation-ladder-certificate-gap}
\label{tab:mechanism-scope-stress-test}
\label{tab:defensive-allocation-comparison}
\label{tab:parameter-sensitivity}
\label{tab:baseline-comparison}
\label{tab:solver-scaling}
\label{tab:gnosis-counterfactual}
\begin{tabularx}{\linewidth}{
@{}
>{\RaggedRight\arraybackslash}p{0.15\linewidth}
>{\RaggedRight\arraybackslash}p{0.30\linewidth}
Y
@{}}
\toprule
Check & Recorded result & Supported inference\\
\midrule
Evidence layers
& Production \(0\); four-of-seven counterfactual
  \(0\to4\to10\)
& Evidence determines the sharp boundary; positive values are conditional.\\

Activation ladder
& Fixed resistance: TCR \(=5\); ACR
  \(=5,14,23,32,41\)
& Ordered evidence excludes cheap static covers.\\

Mechanism stress
& Sequential \(k+4\); package and robust TCR \(=4\)
& A mechanism that bypasses states invalidates activation evidence.\\

Defense allocation
& \(H=10\): local rules \(9\); bottleneck optimum \(11\)
& Effective defense must hit every cheap substitute cover.\\

Interactions
& Lower-order gains \(0\); greedy ratio
  \(\approx5.0\times10^{-5}\)
& Monotonicity gives no singleton-greedy guarantee.\\

Generalization
& Four further shapes, \(3\)-of-\(5\) to \(7\)-of-\(13\), \(400\) seeded
  trials: \(0\) monotonicity failures; greedy exact at budget one throughout
& The certificate and greedy findings are not specific to four-of-seven.\\

Computation
& \(n=18\): \(262{,}144\) states; median \(7.495\) s
& Exact activation search is exponential; uniform TCR uses sorting.\\
\bottomrule
\end{tabularx}
\end{table}

Exact-solver, scaled dynamic-program, sensitivity-grid, baseline, and command
details appear in Appendix~\ref{app:evaluation-artifacts}. Equal-mean
four-of-seven profiles produce certificates from \(16\) to \(40\), confirming
that the lower tail, rather than the mean, drives the uniform certificate.
These checks exercise implementations of stated mathematical inputs; they do
not validate those inputs for real Keypers.

\subsubsection{Activation paths and mechanism scope}

The activation ladder fixes ten equal-weight members, threshold \(t=0.5\),
five cheap members of resistance one, and five seed members of resistance ten.
If \(r\) seed acquisitions are required before the cheap path becomes
available, resistance-only threshold cover remains five while exact activation
cover is
\[
    \operatorname{ACR}_{R,\tau^{(r)}}(\varnothing)
    =
    10r+(5-r)
    =
    5+9r.
\]
For \(r=0,\ldots,4\), this gives
\(5,14,23,32,41\). Because the resistance vector and static threshold cover do
not change, the increase is caused entirely by paths excluded by the certified
activation order.

The mechanism stress test uses \(k\) seed members and two core members.
Sequential acquisition must first buy the seeds and therefore costs \(k+4\);
one-shot independent offers are initially infeasible, while an all-or-nothing
package acquires the core members jointly for four. Mechanism-robust threshold
cover also remains four. The growing sequential--package gap isolates the
mechanism boundary: a package can condition payment on joint completion and
avoid the intermediate exposure states, so simultaneous and package
acquisition are not interchangeable.

\subsubsection{Defensive allocation and interaction order}

The allocation instance has four certified minimal covers over six bottleneck
members and two high-weight decoys. Every minimal cover begins at cost six.
For integer budget \(H\), the bottleneck rule solves the fixed-family max--min
problem, while the two local baselines allocate to the cheapest current
resistance or cycle by nonincreasing member weight.

\begin{table}[!htbp]
\centering
\captionsetup{font=scriptsize,skip=2pt}
\scriptsize
\setlength{\tabcolsep}{6pt}
\renewcommand{\arraystretch}{0.95}
\caption{Certificate under three defensive-allocation rules.}
\label{tab:defense-allocation-detail}
\begin{tabular}{@{}rrrr@{}}
\toprule
Budget \(H\)
& Cheapest resistance
& Weight cycle
& Bottleneck optimum\\
\midrule
0  & 6 & 6 & 6\\
2  & 6 & 6 & 7\\
4  & 6 & 6 & 8\\
6  & 6 & 7 & 9\\
8  & 7 & 9 & 10\\
10 & 9 & 9 & 11\\
\bottomrule
\end{tabular}
\end{table}

At \(H=4\), both local rules leave at least one binding cover unchanged and
return six, whereas the bottleneck allocation reaches eight. At \(H=10\),
both local rules return nine and the optimum returns \(11\). These are
certificate comparisons, not loss estimates: they are valid because the
controlled ledger certifies the outer family.

Exact rational-arithmetic checks cover 18 pure \(k\)-way constructions. Every
nonempty proper M\"obius coefficient is zero, only the full-order coefficient
is nonzero, and omitting it underpredicts the full-defense value by factors up
to \(100\) in the tested family. In a separate seeded study of 100 monotone
four-of-seven certificate functions, singleton greedy is often competitive on
small random instances, but the fixed decoy construction reaches a gain ratio
of approximately \(5.0\times10^{-5}\). Random-instance performance therefore
does not repair the worst-case impossibility result.

\subsubsection{Lower-tail sensitivity and computation}

The sensitivity grid uses four seven-member profiles with identical mean
resistance ten and varies threshold and initial exposure. The left panel
reports the four-of-seven certificate after zero, one, or two initially
exposed shares; the right panel reports ten retained exact-solver runs on the
recorded laptop.

\begin{table}[!htbp]
\centering
\captionsetup{font=scriptsize,skip=2pt}
\fontsize{7}{7.5}\selectfont
\setlength{\tabcolsep}{3pt}
\renewcommand{\arraystretch}{0.95}
\caption{Lower-tail sensitivity and exact-solver scaling.}
\label{tab:sensitivity-scaling-detail}
\begin{subtable}[t]{0.43\linewidth}
\centering
\caption{Four-of-seven certificate.}
\label{tab:lower-tail-panel}
\begin{tabular}{@{}lrrr@{}}
\toprule
Profile & \(e=0\) & \(e=1\) & \(e=2\)\\
\midrule
Uniform & 40 & 30 & 20\\
Balanced & 34 & 24 & 15\\
Bimodal & 16 & 12 & 8\\
Heavy tail & 18 & 10 & 4\\
\bottomrule
\end{tabular}
\end{subtable}
\hfill
\begin{subtable}[t]{0.53\linewidth}
\centering
\caption{Ten retained solver runs.}
\label{tab:exact-solver-panel}
\begin{tabular}{@{}rrrr@{}}
\toprule
\(n\) & States & Median (s) & Peak MiB\\
\midrule
8  & 256     & 0.000683 & 0.0078\\
10 & 1,024   & 0.009757 & 0.0203\\
12 & 4,096   & 0.053806 & 0.0702\\
14 & 16,384  & 0.311079 & 0.2694\\
16 & 65,536  & 1.555494 & 1.0663\\
18 & 262,144 & 7.495092 & 4.2539\\
\bottomrule
\end{tabular}
\end{subtable}
\end{table}

Initial exposure reduces every tested certificate because fewer remaining
shares are required. For the heavy-tail profile at zero exposure, raising the
threshold from three to six shares changes the certificate through
\(10,18,29,45\); all corresponding monotonicities hold over the 48-row
sensitivity grid.

The exact solver follows the \(2^n\) state count
(Table~\ref{tab:exact-solver-panel}); the recorded timing variation is
retained rather than trimmed. These measurements support the exponential
limitation, not large-committee exact scalability -- uniform threshold cover
requires only sorting, while integer-scaled activation inputs may use the
pseudo-polynomial dynamic program.

\subsubsection{Generalization across committee shapes}
\label{subsec:generalized-committee-shapes}

The checks above fix \(n=7\), \(q=4\). A further sweep repeats the seeded
monotonicity, M\"obius truncation-error, and exact-versus-greedy allocation
checks at four additional shapes holding a similar threshold ratio --
\(3\)-of-\(5\), \(5\)-of-\(9\), \(6\)-of-\(11\), and \(7\)-of-\(13\) -- with
\(100\) seeded trials at each. All \(400\) trials are monotone. Median maximum
relative truncation errors after retaining interactions through order one,
two, and three are \(0.310/0.188/0.112\) at \(3\)-of-\(5\),
\(0.341/0.271/0.262\) at \(5\)-of-\(9\), \(0.342/0.251/0.289\) at
\(6\)-of-\(11\), and \(0.374/0.276/0.367\) at \(7\)-of-\(13\) -- comparable in
magnitude to the \(n=7\) values above, and not always decreasing with
retained order (order three exceeds order two at the two larger shapes),
which a signed M\"obius expansion permits and is not a computation error.
Singleton-marginal greedy remains exact at budget one in every shape (minimum
ratio \(1.000\)); at budgets two and three the minimum ratio ranges from
\(0.473\)/\(0.418\) at \(3\)-of-\(5\) to \(0.949\)/\(0.898\) at
\(6\)-of-\(11\), so the earlier finding is not an artifact of \(n=7\).
Truncation error at the larger shapes is computed as a running sum of
per-order zeta transforms rather than the \(O(3^n)\) per-mask enumeration used
above, which keeps the check tractable at \(n=13\); the faster routine matches
the direct enumeration exactly on a held-out instance. As throughout this
subsection, these are seeded controlled instances, not measurements of any
real committee of those sizes.

\subsubsection{Boundary profiles and larger committees}
\label{subsec:boundary-larger-committee}

A second sweep tests the sensitivity grid's open ends and two extreme
profiles: an all-zero resistance vector, and one member holding nearly all
resistance, \((100,1,1,1,1,1,1)\), together with committee sizes \(14\) and
\(28\) built by tiling the original four profiles and holding the
four-of-seven threshold ratio. The all-zero profile certifies exactly zero at
every threshold and exposure combination tested, as it must. The
dominant-member profile holds its certificate at the number of required
shares while unanimity is not required -- \(1\) through \(6\) at
\(q=1,\ldots,6\) -- then jumps to \(106\) once the seventh share, and with it
the dominant member's resistance, becomes unavoidable at \(q=7\): an
eighteenfold increase from the \(q=6\) value, driven entirely by whether the
last share must be acquired. At \(n=14\) (\(q=8\)) the four tiled profiles
certify \(80\), \(68\), \(32\), and \(36\); at \(n=28\) (\(q=16\)) they certify
\(160\), \(136\), \(64\), and \(72\) -- exactly double and quadruple,
respectively, their \(n=7\) values at the matched per-tile share count, as
tiling identical resistances implies. These extend the same normalized
controlled construction to its parameter extremes and introduce no new
production evidence.

\subsubsection{Independent reimplementation}
\label{subsec:independent-reimplementation}

A reimplementation built only from the formulas stated in this paper, not
from the authors' own code, independently cross-checks all 36 numbers named
in this section and stress-tests the central results beyond the hand-picked
examples above. On \(720\) random small instances (\(n=3,\ldots,8\), uniform
and non-uniform weights) plus five hand-picked edge cases -- an all-zero and a
tied activation vector, zero-resistance members, an \(A_0\) just short of the
threshold, and an activation vector admitting no first acquisition -- the
exact activation-cover equivalence (the set-based (AC) reduction against
brute-force sequential search) shows zero mismatches on all \(725\) instances.
The coordinated-hardening and singleton-greedy-no-guarantee constructions are
checked exhaustively -- every proper subset, not a sample -- for \(k\) up to
\(9\) and \(20\) respectively; the activation-boundary instability
construction is confirmed arithmetically for \(K\) up to one million. The
exact-solver scaling table is independently extended to \(n=22\)
(\(4{,}194{,}304\) states) on a separate machine and implementation. The same
reimplementation also checks the replacement-hull characterization of
Theorem~\ref{thm:replacement-hull}: on \(200\) random instances across two
committee/outcome shapes and five values of \(\varepsilon\), the primal and
the \((s,Q)\)-equivalent forms of \(\mathsf A_i^\varepsilon(C,x)\) agree
exactly, checked by extracting \(s^*\) from the primal's own optimal solution
and independently re-solving the equivalent form's inner problem at that
value rather than by a floating-point search; the separation corollary is
confirmed on three constructive instances where the separating functional is
exhibited directly rather than found by a general solver. This
reimplementation catches transcription and implementation errors, not a
conceptual mistake shared by both this paper and the reimplementation.

\subsection{Production Evidence Audit and Controlled Execution}
\label{subsec:public-threshold-cap}
\label{subsec:bonded-slashing-pilot}
\label{subsec:public-input-audit}

We anchor 2026-06-13 00:00 UTC to Gnosis block
\(46{,}666{,}718\), with hash \texttt{0x574e...9b60}. Manager
\texttt{0x7C23...4abb} selects set 10,
\texttt{0xE817...E828}. Archival calls, the creation receipt, and the
\texttt{KeyperSetAdded} event agree on seven ordered members, threshold four,
activation block \(46{,}271{,}880\), and finalized state; the artifact retains
the full addresses, sources, and member-level diagnostics.

The production audit checks direct resistance, activation, attribution,
execution, and insurance or compensation for all seven members. No row contains
a retained auditable positive direct floor, activation floor, or joint lower
bound for attribution, enforcement, and unavoidable loss. A missing entry
means absence of a certified positive floor, not observed zero resistance.
Consequently, the production ledger supports
\(\underline R=(0,0,0,0,0,0,0)\), while actual member resistance remains
unknown. For a four-of-seven committee, a positive threshold-cover certificate
requires at least \(7-4+1=4\) positive member floors, and a target \(B\)
requires the sum of the four smallest certified floors to be at least \(B\).
The present production evidence therefore remains on the public-only zero
branch of Theorem~\ref{thm:information-boundary}.

To show how the nontrivial branches would operate on the same pinned committee
geometry, consider a separate counterfactual ledger
\(\mathcal I_{\mathrm{cf}}\). Keep the seven equal-weight members, threshold
four, and \(A_0=\varnothing\), but suppose future evidence certified, in
normalized cost units,
\[
  (\underline R_i,\underline\tau_i)=
  \begin{cases}
     (4,0),   & i\in\{0,1\},\\
     (1,2/7), & i\in\{2,\ldots,6\}.
  \end{cases}
\]
These values are illustrative assumptions, not measurements or claims about
the current Keypers. Public data alone still certify zero. With only the
counterfactual resistance floors, the four cheapest members give
\(\operatorname{TCR}_{\underline R}(\varnothing)=4\). With the additional
activation floors, members \(2,\ldots,6\) cannot be used before exposure
\(2/7\), so every certified first-crossing path must acquire members \(0\) and
\(1\) before two of the remaining members; hence
\(\operatorname{ACR}_{\underline R,\underline\tau}(\varnothing)=10\).
The same committee geometry therefore produces the sharp boundaries
\(0\to4\to10\) as the counterfactual ledger is strengthened.

The value ten additionally requires a passed activation-scope gate: released
weight must summarize acquisition history, every covered success must retain
an exposure-ordered first-crossing witness, and effective member payments must
be separable after rebates and side transfers. Shared operation, correlated
legal exposure, or any identity-dependent change in resistance would invalidate
that gate and prevent issuance of the activation certificate.

Changing the production result would require dated and scope-matched positive
direct resistance evidence or a penalty route combining member attribution,
an enforceable \(P_i>0\), a positive joint execution floor, and exclusion of
insurance, reimbursement, and double counting. Activation evidence would
further require member-specific state or cap bounds; it can strengthen
activation cover but not mechanism-robust threshold cover.

The separate Chiado execution uses seven validators and seven Keypers on one
host with threshold four. It records seven shares, seven native signatures, a
valid aggregate key, and matching four-share reconstruction. Member \(0\)'s
share was observed \(5{,}350\) seconds before release; an EIP-712 attestation
executed forfeiture in block \(22{,}111{,}234\), reducing the recorded
certificate from \(4\times10^{12}\) to \(3\times10^{12}\) wei after one
\(10^{12}\)-wei loss. The certificate binds nine artifacts, and the live
verifier checks eleven receipts, deployment data, calldata, events, and current
state. Both the pinned production snapshot and the recorded Chiado transaction
trail were independently rechecked against current chain state on 2026-07-20
and returned unchanged.

This execution demonstrates an externally checkable evidence-to-enforcement
transition within the recorded testnet mechanism. It does not establish
independent governance, operation by seven organizations, native on-chain BLS
verification, or production resistance. Production and Chiado are not pooled:
their members, environments, penalties, and evidence procedures differ, and
Dune activity is calibration rather than a worst-case
\(\overline V(x)\). The actual production output therefore remains zero;
issuing activation cover requires separately certified member-level inputs and
a mechanism-matched activation gate.

\subsection{Limitations and Reproducibility}

The audit does not estimate bargaining or actual willingness to collude.
Approximate activation inputs remain unsafe near path boundaries, and general
exact search does not scale. Production evidence also lacks the positive
member floors and worst-case value cap needed for a no-profit claim.

The counterfactual four-of-seven calculation illustrates certificate
computation on the pinned committee geometry, but does not certify its
assumed resistance, activation, operator-independence, or payment conditions
for the production Keypers. The generalized committee-shape sweep and the
boundary and larger-committee sensitivity rows are extensions of the same
seeded or tiled controlled constructions; neither introduces new production or
deployment evidence, and neither claims that a committee of those other sizes
exists. The independent reimplementation catches transcription and
implementation errors but not a conceptual error shared by both this paper
and the reimplementation, and its check of the replacement-hull separation
corollary covers hand-picked constructive instances rather than a general
solver.

The artifact includes the pinned set, seven-row ledger, counterfactual fixture
and result, audit JSON, archival-RPC verifier, Chiado certificate, the
generalized-shape and boundary-sensitivity outputs, and the independent
reimplementation, including its replacement-hull check. Deterministic results
run offline behind a single reproduction command; timing provenance and
live-chain steps appear in Appendices~\ref{app:extended-reproduction}
and~\ref{app:chiado-live-verification}.
